% =============================================================================
% Chapter 2: Philosophy and Design Principles
% =============================================================================

\chapter{Philosophy and Design Principles}
\label{ch:philosophy}

The requirements in this document are derived from a set of foundational
principles. These principles constrain every design decision: an option that
satisfies a requirement but violates a principle is never adopted.

\paragraph{P-01: Data management before analysis.}
A result is only as trustworthy as the data pipeline that produced it.
\itc{} enforces structure, traceability, and explicit intent at every step.
Analysis is downstream of management, never the other way around.

\paragraph{P-02: Provenance and History are mandatory in production.}
Every VarFrame knows where it came from (Provenance) and what was done to it
(History). Both records travel with the data and are preserved in every
export format. Production mode never suppresses either record.

\paragraph{P-03: Reproducibility by design.}
Loading the same files with the same processor and the same equation file
must always produce the same result. There is no implicit state, no silent
default that varies between runs. If the result changes, a hash changes, and
\itc{} reports it.

\paragraph{P-04: Fail fast and loud.}
Ambiguity is an error. Missing variables, grid mismatches, unvalidated
processors, non-monotonic axis-translation targets, mismatched operating
modes, and non-differentiable functions in uncertainty equations raise
informative exceptions before any computation occurs. Silent NaN propagation
and silent fallbacks are unacceptable.

\paragraph{P-05: Test-driven code is non-negotiable.}
\itc{} handles data that feeds engineering decisions and, ultimately, flight
vehicles. Every public method is tested before it is implemented. Coverage
below 90\,\% is a defect, not a metric. Property-based testing is preferred
for components with mathematical contracts (operators, propagation,
smoothing, differentiation, axis rotation).

\paragraph{P-06: Immutability by default.}
Operations on a VarFrame return new objects. The original is never modified.
This makes reasoning about data pipelines straightforward and eliminates a
large class of bugs. Structural immutability is enforced mechanically, not by
convention.

\paragraph{P-07: Minimal API surface.}
No method is added if it can be composed from existing ones. Every public
function traces back to a concrete, documented use case. Complexity is the
enemy of reliability.

\paragraph{P-08: Designed for notebooks.}
Inspection, summary, and diagnostic methods print rich output to the terminal
and return inspectable objects. The workflow must feel natural in Jupyter
without additional boilerplate. No method exists that only prints or only
returns: every diagnostic surface offers both faces.

\paragraph{P-09: English everywhere.}
All code, comments, docstrings, variable names, error messages, equation
files, and documentation must be written in English. American English with
Z is used consistently (\emph{organize}, \emph{visualization},
\emph{normalize}). Development conversation may occur in any language; the
artifacts must not.

\paragraph{P-10: Example-driven development.}
Every feature is specified as a usage example before implementation begins.
The implementation serves the example, not the reverse. If a feature cannot
be demonstrated in a clear usage example, it does not exist.

\paragraph{P-11: Standards alignment is intentional, not coincidental.}
\itc{} explicitly aligns with established standards: the GUM for
uncertainty, AIAA for coordinate systems and wind tunnel uncertainty
assessment, SAE for in-flight thrust determination, W3C PROV-DM for the
conceptual provenance model. Where \itc{} departs from a standard, the
departure is documented and justified in
Chapter~\ref{ch:standards}.

\paragraph{P-12: Sensitivity is a native capability.}
Sensitivity analysis is not a separate module. It is the direct output of
combining \code{set\_uncertainty} on input variables with \code{compute} or
\code{itc.aerospace} computations: the propagated uncertainty on the output
quantifies its sensitivity to the specified input variation. This is a
deliberate architectural property, not a coincidence.

\paragraph{P-13: Explicit over implicit numerical combination.}
Arithmetic combinations of two VarFrames go through \code{db.combine(other,
op=...)}, not through overloaded Python operators. This forces every
numerical combination to record its semantics in the History and to declare
how uncertainties and correlations are handled. Convenience that hides
numerical assumptions is not convenience; it is a defect.

\paragraph{P-14: Memory is not free.}
Operations are vectorized via NumPy and avoid unnecessary copies whenever
correctness allows. Operations that may inflate memory (\code{expand},
broadcasts, dense pivots from sparse acquisition matrices) document their
expected memory footprint and offer chunked or sparse-aware alternatives
where appropriate. Large or sparse datasets are first-class citizens, not
edge cases.
